\documentclass[11pt]{article}
\usepackage[a4paper,margin=24mm]{geometry}
\usepackage{amsmath,amssymb,amsthm,mathtools}
\usepackage{booktabs,array,longtable,enumitem}
\usepackage[hidelinks]{hyperref}
\usepackage{microtype}

\newtheorem{theorem}{Theorem}[section]
\newtheorem{lemma}[theorem]{Lemma}
\newtheorem{corollary}[theorem]{Corollary}
\theoremstyle{definition}
\newtheorem{definition}[theorem]{Definition}
\newtheorem*{reviewquestion}{Reviewer question}

\newcommand{\TT}{\mathrm{TT}}
\newcommand{\rank}{\operatorname{rank}}
\newcommand{\Mat}{\operatorname{Mat}}
\newcommand{\CDet}{\operatorname{CDet}}
\newcommand{\perm}{\operatorname{perm}}
\newcommand{\FEMPS}{\textnormal{\textsc{femps}}}
\newcolumntype{P}[1]{>{\raggedright\arraybackslash}p{#1}}
\newcommand{\reviewline}{\par\smallskip\noindent
  \textbf{Reviewer decision:} $\square$ verified \quad $\square$ revision
  needed \quad $\square$ outside expertise\par\medskip}

\title{Evidence Audit Companion to\\
\emph{Fermionic Extensions of Continuous Functional Tensor Networks}}
\author{Prepared for independent human review}
\date{}

\begin{document}
\maketitle

\begin{abstract}
This companion reorganizes the evidence behind the associated PRA manuscript
for line-by-line human review.  It is not a second paper and introduces no
additional scientific claim.  For every major conclusion it states the input
domain, the proof dependency, the reduction or calculation, the quantities
that can be independently recomputed, and the boundary beyond which the
conclusion does not extend.  The final section gives a complete audit trail for
the selected interacting calculation, including the state class, Hamiltonian,
operator construction, optimization schedule, initialization lineage,
reference calculation, error measures, resource measurements, and
antisymmetry residual.  Finite-size arithmetic checks can detect transcription
errors, but only the displayed arguments and the cited source theorems support
claims at arbitrary size.
\end{abstract}

\tableofcontents

\section{How to use this companion}

The associated PRA manuscript is the publication text.  This document is a
review aid.  A reviewer may audit the mathematical claims independently of the
floating-point calculation and should record any unverified dependency rather
than infer it from a successful small case.  In particular, the two reductions
that use Chien--Harsha--Sinclair--Srinivasan (CHSS)
\cite{ChienHarshaSinclairSrinivasan2011NoncommDet} require independent checking
of the cited source theorem, its field assumptions, and the row-order
convention.

The following boundaries apply throughout.
\begin{enumerate}[leftmargin=*,itemsep=2pt]
\item Exact worst-case hardness does not exclude Monte Carlo, additive
  approximation, or an efficiently contractible promised subclass.
\item The ordinary particle-TT rank statements do not apply automatically to
  occupation-number MPS, ordered-coordinate methods, or exterior cores.
\item A floating-point result is numerical evidence, not a theorem.
\item The selected numerical state is a finite nonorthogonal Slater expansion
  and is therefore not claimed to be distinct from NOCI.
\item Antisymmetry is exact in the exterior representation; every numerical
  summary nevertheless reports the corresponding residual.
\end{enumerate}

\section{Claim map and dependency structure}

\begin{longtable}{@{}P{15mm}P{45mm}P{50mm}P{30mm}@{}}
\toprule
Item & Claim & Principal support & Independent dependency \\
\midrule
\endfirsthead
\toprule
Item & Claim & Principal support & Independent dependency \\
\midrule
\endhead
E1 & Exact particle-TT bonds equal particle-cut matrix ranks. & Successive
rank factorizations (TT-SVD). & Oseledets \cite{Oseledets2011TT}. \\
E2 & Every nonzero alternating $N$-tensor has cut rank at least
$\binom Nk$. & Explicit diagonal minor selected from one nonzero Pluecker
coordinate. & None. \\
E3 & A Slater determinant has a flat particle Schmidt spectrum and the stated
optimal rank-$r$ error. & Partition of the permutation sum plus
Eckart--Young. & Coleman \cite{Coleman1963FermionRDM}. \\
E4 & Exact squared norm is \#P-hard for rational one-form \FEMPS{} at maximum
bond two. & Direct Cayley-coefficient lemma and one-query metric reduction. &
CHSS Theorems 3.5 and 3.9. \\
E5 & Exact point evaluation is \#P-hard in an explicit rational shifted-
Legendre basis at maximum bond two. & CHSS input, invertible evaluation
matrix, explicit bit bounds, and one-query postprocessing. & CHSS plus standard
Vandermonde identities. \\
E6 & Specified matrix-pair memories collapse to polynomial LC--AGP families. &
Radical filtration or fixed-counter interpolation followed by polarization. &
Supplied algebra decomposition or embedding is part of the input. \\
E7 & A bandwidth-one pair path encodes a permanent; real-PSD relative norm
inherits the stated obstruction. & Unique path and even-form commutativity. &
Valiant \cite{Valiant1979Permanent}; Meiburg \cite{Meiburg2023PSDPermanent}. \\
E8 & A universal direct statistics-carrier tensor product fails. & Explicit
two-Slater form with one-cut rank $N+2$. & None. \\
E9 & The selected $N=6,D=12,K=4$ interacting calculation has the reported
energy, error, variance, norm, cost, and exact antisymmetry. & Determinant
transition solver and independent same-basis exterior CI. & Numerical
recomputation; no theorem inference. \\
\bottomrule
\end{longtable}

\section{Common definitions and normalization}

Let $V_D$ be a $D$-dimensional real or complex one-particle space with
orthonormal basis $e_1,\ldots,e_D$.  We use the isometric alternating embedding
\begin{equation}
J_N(v_1\wedge\cdots\wedge v_N)=\frac{1}{\sqrt{N!}}
\sum_{\pi\in S_N}\operatorname{sgn}(\pi)
v_{\pi(1)}\otimes\cdots\otimes v_{\pi(N)}.
\label{audit:isometry}
\end{equation}
A one-form matrix-wedge state has cores
$A^{[j]}\in\Mat_{\chi_{j-1}\times\chi_j}(V_D)$,
$\chi_0=\chi_N=1$, and
\begin{equation}
C(A)=\bigl[A^{[1]}\wedge\cdots\wedge A^{[N]}\bigr]_{11}
\in\Lambda^N V_D.
\label{audit:femps}
\end{equation}
Every virtual path is a wedge of $N$ one-forms.  Thus Eq.~\eqref{audit:femps}
is alternating without an approximate projection.  The structural
antisymmetry residual is therefore exactly zero.  Whenever a labelled-particle
tensor is materialized for a small system, an independent residual is
\begin{equation}
\epsilon_{\rm AS}^{\rm mat}
=\max_{1\le i<N}
\frac{\|C+\tau_{i,i+1}C\|_2}{\|C\|_2},
\label{audit:asres}
\end{equation}
where $\tau_{i,i+1}$ swaps adjacent particle slots.

\reviewline

\section{E1--E3: exchange rank in particle factorization}

\subsection{E1: exact particle-TT rank}

\begin{theorem}
For every coefficient tensor $C$ and cut $k\mid N-k$, the smallest exact
ordinary particle-TT bond is
\begin{equation}
r_k^{\TT}(C)=\rank C_{(k)}.
\end{equation}
\end{theorem}

\begin{proof}
Contract the first $k$ cores of any exact TT into a matrix $L_k$ and the
remaining cores into $R_k$.  Then $C_{(k)}=L_kR_k$, so the bond is at least
$\rank C_{(k)}$.  Conversely, factor the first unfolding at its exact rank,
reshape the right factor, and repeat.  At the $k$th step the carried index has
dimension exactly $\rank C_{(k)}$.  The successive factorizations reconstruct
$C$, establishing simultaneous attainability at all cuts.
\end{proof}

\begin{reviewquestion}
Does the proof establish simultaneous TT ranks, rather than unrelated
factorizations at individual cuts?  Is the theorem restricted to the
labelled-particle tensor product?
\end{reviewquestion}
\reviewline

\subsection{E2: universal exchange floor}

\begin{theorem}
If $0\ne C\in\Lambda^N V_D$, then
$\rank C_{(k)}\ge\binom Nk$ for $1\le k<N$.  Equality is attained by every
nonzero decomposable $N$-form.
\end{theorem}

\begin{proof}
Choose an $N$-element set $I$ whose alternating coefficient $c_I$ is nonzero.
Restrict the row index of $C_{(k)}$ to the $k$-subsets $A\subset I$ and the
column index to their complements $I\setminus A$.  The resulting
$\binom Nk\times\binom Nk$ submatrix is diagonal up to row and column signs:
the matched entry is $\pm c_I$, while a mismatched row and column repeat at
least one basis index and hence have zero alternating coefficient.  Its
determinant is nonzero, proving the lower bound.  If
$C=u_1\wedge\cdots\wedge u_N$, all contractions lie in
$\Lambda^{N-k}\operatorname{span}\{u_1,\ldots,u_N\}$, which has dimension
$\binom Nk$; the preceding lower bound is therefore sharp.
\end{proof}

The immediate truncation consequence is that a nonzero ordinary particle TT
with a smaller bond at any cut cannot remain exactly alternating.  This does
not say that a mode-space MPS needs the same bond.

\reviewline

\subsection{E3: flat Slater spectrum}

\begin{theorem}
For orthonormal $u_1,\ldots,u_N$, the normalized Slater state $\Phi$ has
\begin{equation}
\Phi=\binom Nk^{-1/2}\sum_{|I|=k}(-1)^{\epsilon(I)}
\Phi_I\otimes\Phi_{I^c}.
\label{audit:slater-schmidt}
\end{equation}
All $\binom Nk$ singular values are $\binom Nk^{-1/2}$, and the best cut-rank
$r$ approximation has relative squared error
$1-r/\binom Nk$.
\end{theorem}

\begin{proof}
Partition the $N!$ signed permutations in Eq.~\eqref{audit:isometry} by the set
$I$ of orbitals appearing in the first $k$ slots.  Each class contains
$k!(N-k)!$ terms and factors into normalized Slater states $\Phi_I$ and
$\Phi_{I^c}$ with coefficient
$\sqrt{k!(N-k)!/N!}=\binom Nk^{-1/2}$.  Different subsets are orthogonal on
both sides.  Equation~\eqref{audit:slater-schmidt} is therefore a Schmidt
decomposition.  Discarding the smallest singular values is optimal by the
Eckart--Young theorem; because they are all equal, retaining $r$ leaves the
stated squared error.
\end{proof}

\begin{reviewquestion}
Confirm that the orbitals are orthonormal and that no flat-spectrum statement
is made for a general correlated alternating tensor.
\end{reviewquestion}
\reviewline

\section{E4: fixed-bond exact squared-norm obstruction}

For a matrix $M=(M_{ij})$ over a noncommutative algebra, define the
row-ordered Cayley determinant
\begin{equation}
\CDet(M)=\sum_{\sigma\in S_N}\operatorname{sgn}(\sigma)
M_{1\sigma(1)}\cdots M_{N\sigma(N)}.
\label{audit:cayley}
\end{equation}
The product order is increasing in the row label.  This is the essential
difference from the ordinary scalar determinant.

\begin{lemma}[direct coefficient]
Set $\mathcal A^{[i]}=\sum_jM_{ij}e_j$.  Then for boundary vectors $u,v$,
\begin{equation}
u^{\mathsf T}\mathcal A^{[1]}\wedge\cdots\wedge
\mathcal A^{[N]}v
=\bigl(u^{\mathsf T}\CDet(M)v\bigr)e_1\wedge\cdots\wedge e_N.
\label{audit:direct}
\end{equation}
\end{lemma}

\begin{proof}
A top-form term must select each $e_j$ exactly once, hence is indexed by a
permutation.  Reordering its exterior factors contributes the permutation
sign, while matrix multiplication retains row order.  Summation gives
Eq.~\eqref{audit:direct}.
\end{proof}

\begin{theorem}
Over $\mathbb Q$, exact squared-norm evaluation for rational one-form
\FEMPS{} is \#P-hard under a polynomial-time metric reduction even when
$D=N$ and all internal bonds are at most two.
\end{theorem}

\begin{proof}
For a 3CNF formula $\varphi$ with $m$ clauses and $S$ satisfying assignments,
CHSS Theorems 3.5 and 3.9 construct in polynomial time a polynomial-size
matrix $H_\varphi$ over $\Mat_2(\mathbb Q)$ with constant-bit integer gadget
weights and
\begin{equation}
\CDet(H_\varphi)=aI_2+bJ_2,
\qquad a+b=4^{3m}S,
\qquad J_2=\begin{pmatrix}0&1\\1&0\end{pmatrix}.
\label{audit:chss-output}
\end{equation}
Use Eq.~\eqref{audit:direct} with $u=e_1$ and $v=e_1+e_2$.  The state has only
one top-form coefficient,
$u^{\mathsf T}(aI_2+bJ_2)v=a+b=4^{3m}S\ge0$, and therefore norm
$16^{3m}S^2$.  A single exact norm query, exact nonnegative integer square
root, and division by $4^{3m}$ recover $S$.  The CHSS matrix size and core
encoding are polynomial in $|\varphi|$; the queried integer has polynomial bit
length because $S\le2^{n_{\rm var}}$.  Both the construction and
postprocessing are polynomial-time, and the maximum bond stays two.
\end{proof}

\paragraph{External dependency to check.}
The proof uses the structured CHSS output in Eq.~\eqref{audit:chss-output}, not
merely hardness of an arbitrary signed matrix entry.  The reviewer should
confirm: (i) CHSS uses the same row-ordered convention; (ii) the construction
is valid over characteristic zero; (iii) its gadget weights and graph size
have polynomial encoding; and (iv) the stated boundary indeed selects $a+b$.
For a generic signed Cayley entry the manuscript retains only the weaker
bond-three polarization observation.

\reviewline

\section{E5: exact point evaluation in a rational Legendre basis}

Let $\ell_r(t)=P_r(2t-1)$ be the unnormalized shifted Legendre polynomial,
$P_r(1)=1$.  Its integer expansion is
\begin{equation}
\ell_r(t)=\sum_{k=0}^r(-1)^{r-k}\binom rk\binom{r+k}{k}t^k,
\label{audit:legendre}
\end{equation}
with leading coefficient $\binom{2r}{r}$.

\begin{theorem}
Exact point evaluation of rational one-form \FEMPS{} in
$\ell_0,\ldots,\ell_{D-1}$ is \#P-hard under a polynomial-time metric
reduction already at $D=N$ and maximum bond two.  Under the isometry
Eq.~\eqref{audit:isometry}, the output is represented exactly as
$q/\sqrt{N!}$ with rational $q$.
\end{theorem}

\begin{proof}
Choose rational nodes $t_i=i/(N+1)$ and set
$B_{ij}=\ell_{j-1}(t_i)$.  Degree grading changes the monomial Vandermonde
matrix by an upper-triangular coefficient matrix, so
\begin{equation}
\det B=\left(\prod_{r=0}^{N-1}\binom{2r}{r}\right)
\prod_{1\le i<k\le N}(t_k-t_i)\ne0.
\label{audit:Bdet}
\end{equation}
The nodes have $O(\log N)$ bits.  Equation~\eqref{audit:legendre} gives
polynomial-bit evaluations.  More explicitly,
$Q=(N+1)^{N-1}$ clears all denominators in $B$.  The integer entries of $QB$
have polynomial bit length; Hadamard bounds applied to its determinant and
minors imply polynomial bit length for
\begin{equation}
B^{-1}=Q\,\operatorname{adj}(QB)/\det(QB).
\end{equation}
Thus the basis evaluation map and its inverse are exactly computable in
polynomial Turing time.

Take the CHSS matrix $H_\varphi$ from E4 and define
\begin{equation}
\mathcal A^{[i]}(t)=\sum_{j=1}^N H_{ij}\ell_{j-1}(t),
\qquad u=e_1,\quad v=e_1+e_2.
\end{equation}
In the alternating point expansion at $(t_1,\ldots,t_N)$, repeated basis
columns cancel.  Each column permutation contributes its sign times $\det B$
and leaves the CHSS matrix product in row order.  Hence
\begin{equation}
\Psi(t_1,\ldots,t_N)
=\frac{\det B}{\sqrt{N!}},4^{3m}\#\mathrm{SAT}(\varphi).
\label{audit:point}
\end{equation}
The exact algebraic output is the pair $(q,N)$ representing $q/\sqrt{N!}$;
$N!$ has $O(N\log N)$ bits.  One query followed by rational division by the
known nonzero factor $4^{3m}\det B$ recovers $\#\mathrm{SAT}(\varphi)$.
\end{proof}

\paragraph{Separation from E4.}
E4 is a norm query on a coefficient state; E5 is a point query in a specified
functional basis.  They are separate reductions from the same CHSS source.
Neither implies that all $\chi\ge2$ instances are unsuitable for QMC or VMC.

\reviewline

\section{E6--E7: tractable corridors and a sparse hard path}

\subsection{Bounded coefficient algebra}

Let $\mathcal A=S\oplus R$ be a finite-dimensional complex algebra with
$R^d=0$ and
$S=\bigoplus_{a=1}^q\Mat_{p_a}(\mathbb C)$; put
$p=\max p_a$ and $\rho=\dim R$.  Expanding $\Omega^M$ by the number $k$ of
radical factors leaves only $k<d$.  Resolving the $k+1$ semisimple runs into
Wedderburn blocks and the radical insertions in a basis gives at most
$q^{k+1}\rho^k$ structural choices.  For each choice, even physical two-forms
commute, and after the boundary functional is applied one obtains a homogeneous
degree-$M$ polynomial in at most
\begin{equation}
v_k=(k+1)p^2+k
\end{equation}
commuting two-forms.  Degree-$M$ powers span this polynomial space in
characteristic zero.  Therefore
\begin{equation}
K\le\sum_{k=0}^{\min(d-1,M)}q^{k+1}\rho^k
\binom{M+v_k-1}{v_k-1}
\end{equation}
AGP terms suffice.  For the rational Turing statement, the displayed
Wedderburn--Malcev decomposition and radical basis must be supplied; the
theorem does not claim an algorithm that discovers them.

\subsection{Fixed-state graded memory}

If the coefficient algebra is supplied with an embedding in
\begin{equation}
\Mat_w\!\left(\mathbb C[z_1,\ldots,z_g]/
(z_1^{d_1},\ldots,z_g^{d_g})\right)
\end{equation}
for fixed $w,g$, lift the quotient entries and note
$\deg_{z_j}\Omega^M\le M(d_j-1)$.  Tensor-product Vandermonde interpolation
extracts the counter coefficients.  At each grid point the boundary of the
matrix power is a homogeneous polynomial in at most $w^2$ commuting
two-forms, so polarization yields
\begin{equation}
K\le\prod_{j=1}^g[M(d_j-1)+1]
\binom{M+w^2-1}{w^2-1}.
\end{equation}
Rational nodes, Hadamard bounds, and exact Gaussian elimination give the
polynomial bit bound when the embedding is part of the input.

\subsection{Sparse APG permanent path}

For orthonormal pair forms $P_j=e_{2j-1}\wedge e_{2j}$ and a zero--one matrix
$A$, set $F_i=\sum_jA_{ij}P_j$.  The unique endpoint path of an upper-
bidiagonal pair matrix produces
\begin{equation}
\Psi_A=\frac{F_1\wedge\cdots\wedge F_M}{M!}
=\frac{\perm(A)}{M!}P_1\wedge\cdots\wedge P_M.
\end{equation}
Repeated $P_j$ vanish, and distinct even pair forms commute, so precisely the
permanent terms survive.  The squared norm is
$\perm(A)^2/(M!)^2$.  Exact nonnegative square root recovers the zero--one
permanent.  For real positive-semidefinite rational $A$, a relative squared-
norm approximation would similarly give a relative approximation to
$\perm(A)$; the manuscript uses Meiburg's result for the corresponding
complexity consequence.  Entrywise-nonnegative permanents have an FPRAS
\cite{JerrumSinclairVigoda2004PermanentFPRAS}, so the promises must not be
conflated.

\reviewline

\section{E8: failure of a universal direct statistics carrier}

For $C\in\Lambda^N V$, consider the one-cut contraction map
$c_{C,1}:V^*\to\Lambda^{N-1}V$.  A universal direct factorization with a
Slater-normalized statistics carrier would force every admitted one-cut rank
to be divisible by $N$, because a Slater has rank $N$ and multiplicity one.

For $N\ge3$ define in $V=\mathbb F^{N+2}$
\begin{equation}
C_N=e_1\wedge(e_2\wedge e_3+e_4\wedge e_5)
\wedge e_6\wedge\cdots\wedge e_{N+2}.
\end{equation}
This is a sum of two Slaters.  Contracting successively by the duals of
$e_1,\ldots,e_{N+2}$ gives $N+2$ forms with distinct monomial supports, hence
linearly independent images.  Therefore $\rank c_{C_N,1}=N+2$, while
$N\nmid N+2$.  The proposed universal direct tensor product cannot hold for
all nonzero $N$-forms.  The proof does not exclude state-adaptive
factorizations, Hamiltonian-specific charges, or a different categorical
carrier with its own reconstruction theorem.

\reviewline

\section{E9: selected interacting numerical calculation}

\subsection{State class and physical model}

The numerical state is the nonbranching diagonal-path subclass
\begin{equation}
\Psi=\sum_{a=1}^{K}c_a\Phi_a,
\qquad
\Phi_a=u_{a1}\wedge\cdots\wedge u_{aN}.
\label{audit:noci}
\end{equation}
Equation~\eqref{audit:noci} is exactly a $K$-term nonorthogonal Slater
expansion.  It is a restricted \FEMPS{} realization, but not a state class
beyond NOCI.  The selected calculation uses $N=6$, $D=12$, and $K=4$ in the
first twelve one-dimensional harmonic-oscillator functions.  The Hamiltonian
is
\begin{equation}
H=\sum_{i=1}^{N}\left(-\frac12\partial_{x_i}^2+\frac12x_i^2\right)
+\sum_{i<j}\frac{1}{\sqrt{(x_i-x_j)^2+1}}.
\label{audit:soft}
\end{equation}

The soft-Coulomb tensor is projected with 128-point Gauss--Hermite quadrature.
A 160-point calculation changes that tensor relatively by
$1.3545\times10^{-12}$.  A physical-index SVD retains 23 factors at $D=12$;
the relative reconstruction error is $3.8469\times10^{-15}$.  The independent
reference uses the unfactorized four-index quadrature tensor and direct
Slater--Condon diagonalization in the
$\binom{12}{6}=924$ dimensional exterior space.

\subsection{Determinant transition algorithm}

Let $U_a\in\mathbb C^{D\times N}$ collect the orbitals in $\Phi_a$.  For each
of the $K^2$ ordered transition pairs,
\begin{equation}
S_{ab}=U_a^\dagger U_b,
\qquad \langle\Phi_a|\Phi_b\rangle=\det S_{ab}.
\end{equation}
For a one-body matrix $h$, with $X_{ab}=U_a^\dagger hU_b$,
\begin{equation}
\langle\Phi_a|H_1|\Phi_b\rangle
=\det(S_{ab})\operatorname{Tr}(S_{ab}^{-1}X_{ab}).
\label{audit:onebody}
\end{equation}
For a factorized two-body term $L_\ell\otimes R_\ell$, set
$A_{ab}^{\ell}=U_a^\dagger L_\ell U_b$ and
$B_{ab}^{\ell}=U_a^\dagger R_\ell U_b$.  Its transition value is
\begin{equation}
\det(S_{ab})\left[
\operatorname{Tr}(S_{ab}^{-1}A_{ab}^{\ell})
\operatorname{Tr}(S_{ab}^{-1}B_{ab}^{\ell})
-\operatorname{Tr}(S_{ab}^{-1}A_{ab}^{\ell}
S_{ab}^{-1}B_{ab}^{\ell})\right].
\label{audit:twobody}
\end{equation}
Near a singular overlap, the implementation uses column-replacement cofactors
instead of an inverse.  In the selected run all 16 transition pairs passed the
conditioning criterion and used the inverse route.  No virtual path history
and no $D^N$ labelled-particle coefficient tensor were enumerated.

\subsection{Optimization and initialization lineage}

Each $U_a$ is returned to a Stiefel gauge by reduced QR with deterministic
diagonal phases.  At every orbital step the amplitudes are eliminated by the
conditioned generalized Hermitian eigenproblem
\begin{equation}
H(U)c=E\,S(U)c.
\end{equation}
All four overlap directions were retained in the selected run; the final
condition number was 2.1111 and the generalized residual norm was
$1.96\times10^{-14}$.  Automatic differentiation propagates through QR,
transition contractions, and the generalized eigenvalue.

The nonlinear schedule used 160 Adam steps with a cosine learning-rate change
from $10^{-3}$ to $10^{-5}$, followed by 80 requested L-BFGS refinement steps
at learning rate 0.5.  The random seed was 3212.  The starting orbitals were
not a fresh $D=12$ Slater and were not derived from CI.  They were obtained by
exact zero-padding of an already optimized $D=10,K=4$ state.  The energy moved
from 25.0498252875 at the start to 25.0494882497 after Adam and to
25.0494711446 after L-BFGS.  This lineage must remain visible because it
precludes a claim of blind-from-scratch $D=12$ optimization.

\subsection{Definitions of the reported checks}

For the normalized exterior coefficient vector $c_{\rm ext}$ and the dense
same-basis Hamiltonian $H_{\rm CI}$,
\begin{align}
n&=\langle c_{\rm ext},c_{\rm ext}\rangle,\\
E&=\langle c_{\rm ext},H_{\rm CI}c_{\rm ext}\rangle/n,\\
\sigma_E^2&=\|(H_{\rm CI}-E)c_{\rm ext}\|_2^2/n.
\end{align}
The energy error is $E-E_{\rm CI}$, where $E_{\rm CI}$ is the lowest
eigenvalue of the same 924-dimensional dense exterior Hamiltonian.  The norm
error is $|n-1|$.  Structural antisymmetry is zero by Eq.~\eqref{audit:noci};
a full $12^6$ particle tensor was deliberately not built, so no materialized
residual is substituted for the structural result.

\subsection{Results and resource account}

\begin{table}[ht]
\centering
\begin{tabular}{@{}lr@{}}
\toprule
Quantity & Value \\
\midrule
Variational energy & 25.0494711446 \\
Same-basis CI energy & 25.0493664161 \\
Error relative to CI & $1.0473\times10^{-4}$ \\
Energy variance & $1.1223\times10^{-3}$ \\
Norm error & $1.55\times10^{-15}$ \\
Structural antisymmetry residual & 0 \\
Factorized/dense energy difference & $2.84\times10^{-14}$ \\
Wall time & 107.7 s on CPU \\
Peak sampled process RSS & 974,974,976 bytes \\
Exterior dimension & 924 \\
Transition pairs & 16 \\
Stored orbital and amplitude scalars & 292 \\
Enumerated virtual paths & 0 \\
Materialized particle coefficients & 0 \\
\bottomrule
\end{tabular}
\caption{Selected numerical result.  The peak RSS is a sampled whole-process
measurement, not an asymptotic memory bound.}
\label{audit:numtable}
\end{table}

The absolute FEMPS/NOCI and direct-CI energies both decrease along the tested
$D=8,10,12$ basis sequence, but the error relative to same-basis CI is not
monotone at fixed $K=4$.  Consequently Table~\ref{audit:numtable} is a
finite-basis numerical illustration, not evidence for accelerated $D$
convergence.  It establishes neither a new ansatz beyond NOCI nor superiority
over Li--Waintal, DMRG, particle TT, or FCI.

\begin{reviewquestion}
Recompute the dense exterior norm, Rayleigh quotient, residual variance, and
lowest CI eigenvalue from the final 924 coefficients.  Confirm that the
factorized and unfactorized Hamiltonian energies agree within the stated
tolerance and that the zero-padded initialization is disclosed in any summary
of the result.
\end{reviewquestion}
\reviewline

\section{Final human sign-off}

The following sign-off is intentionally not prefilled.

\begin{tabular}{@{}P{45mm}P{95mm}@{}}
Reviewer name and affiliation: & \\[8mm]
Area of expertise: & \\[8mm]
Items reviewed (E1--E9): & \\[8mm]
Unresolved objections: & \\[16mm]
Required revisions: & \\[16mm]
Overall recommendation: & $\square$ accept evidence as stated \quad
$\square$ revise \quad $\square$ reject \\[8mm]
Signature/date: & \\
\end{tabular}

\bibliographystyle{unsrt}
\bibliography{references/references}

\end{document}
